\documentclass[11pt]{article}
\usepackage[a4paper,margin=21mm]{geometry}
\usepackage[no-math]{fontspec}
\setmainfont{Latin Modern Roman}
\newfontfamily\CJKfont{Noto Serif CJK SC}
\XeTeXlinebreaklocale "zh"
\XeTeXlinebreakskip=0pt plus 1pt
\usepackage{amsmath,amssymb,amsthm,amscd,graphicx,color}
\usepackage{xurl}
\usepackage[unicode,hidelinks]{hyperref}
\allowdisplaybreaks
\setlength\emergencystretch{3em}
\numberwithin{equation}{section}

\usepackage[all]{xy}

\numberwithin{equation}{section}

\newtheorem{Theorem}{Theorem}[section]
\newtheorem*{Theorem*}{Theorem}
\newtheorem{Corollary}[Theorem]{Corollary}
\newtheorem{Lemma}[Theorem]{Lemma}
\newtheorem{Proposition}[Theorem]{Proposition}
 { \theoremstyle{definition}
\newtheorem{Definition}[Theorem]{Definition}
\newtheorem{Note}[Theorem]{Note}
\newtheorem{Example}[Theorem]{Example}
\newtheorem{Remark}[Theorem]{Remark} }
\newtheorem{Conjecture}[Theorem]{Conjecture}


\makeatletter\newlabel{section3}{{3}{}{Original source locator}{}{}}\makeatother
\makeatletter\newlabel{prop2.2}{{3.3}{}{Original source locator}{}{}}\makeatother
\begin{document}
\begin{center}\Large Generalized p-Brauer supercharacters of cyclic Tate cohomology are explicit trace combinations for finite free p-adic lattices with central composite-order action\end{center}
\noindent\textbf{Stable ID:} 1473\quad\textbf{Author:} Satoru Urano\par
\noindent\textbf{Work:} A Composite Order Generalization of Modular Moonshine\par
\noindent\textbf{Source:} \url{https://sigma-journal.com/2021/110/}\par
\noindent\textbf{Version:} Official SIGMA 17 (2021) article 110, DOI 10.3842/SIGMA.2021.110; journal PDF and native TeX acquired 2026-09-30, exact bytes pinned by SHA256\par
\noindent\textbf{Rights:} CC BY 4.0. Authors retain copyright. \url{https://creativecommons.org/licenses/by/4.0/}. Reuse and typesetting adaptation; original language and mathematical content preserved.\par
\noindent\textbf{Difficulty (prior selection preserved):} {\CJKfont H2: The proof handles finite-length modules over ramified valuation rings rather than ordinary modular vector spaces. It proves ramification-invariance of the normalized Brauer character, computes cyclic Tate cohomology lengths for eigenlattices, and uses inverse finite Fourier transforms and cyclotomic sums to obtain the exact trace coefficients. These length and ramification obligations go beyond the prime-order Brauer character formula.}\par
\medskip\noindent\textbf{Editorial boundary note (not author text):} Selected author-main Theorem5.7 super p-Brauer trace formula in the fixed base-ring scope R=Z\_p, p dividing N, for finite free integral lattices with central composite-order g and N-regular h. Full regularity/Brauer/Tate/cyclic parity/exact sequence/ramification/Frobenius/discrete Fourier conventions, complete Tate flat base change and rank-one eigenlattice calculation, normalized generalized-character definition/choice/additivity/ramification-invariance, full three-case cyclotomic length calculation and exact conjugated Fourier/Mobius coefficient proof are retained. Original broader R wording remains literal, while the selected integral-lattice scope and actual finite cyclotomic base change are explicit. Standard primitive-root valuation and finite-field/Fourier facts remain precise references. Original Lemma5.6 last scalar-length display suppresses the h-eigenvalue factor while its full statement/opening paragraph retain it; all expressions, conjugation bar and source index variants are preserved without formula repair. Monster vanishing/counterexample/computational consequences are unselected and supporting units have no extra counts. All mathematics and XY diagrams remain native editable TeX; no editor proof.\par
\noindent\textbf{External original reference section3:} Original unselected prime-order modular moonshine history Section3/native212–318; introduction mentions it, but no Monster/vertex-algebra vanishing result is a selected-proof premise.\par
\noindent\textbf{External original reference prop2.2:} Original prime-order Proposition3.3/native240–245, attributed to Borcherds–Ryba1996 Proposition2.2; stated only as the prior case being generalized. Full new composite-order proof is retained.\par
\medskip\hrule\medskip\noindent\textbf{Author text begins below.}\par
\setcounter{section}{1}
% Original source lines 61--63
It is natural to consider Tate cohomology of~$V$ for elements of composite order, but there are no results in the literature. Tate cohomology does not yield vector spaces over a field when the element has composite order, but Brauer characters are only defined for such objects. Specifically, Tate cohomology yields torsion modules over the ground ring. We can't use the Brauer character for such modules because it is defined only for representations defined over a~field and $p$-regular elements. So, we need to define a generalized Brauer character.

In Section~\ref{section2}, we introduce some terminology in this paper. In Section~\ref{section3}, we summarize the modular moonshine conjecture for elements of prime order. In Section~\ref{section4}, we consider ramified extensions $R_p$ of the $p$-adic integers $\mathbb{Z}_p$ and calculate Tate cohomology. In Section~\ref{section5}, to generalize the modular moonshine conjectures, we define a generalized Brauer character and show that the generalized super Brauer character of Tate cohomology is a linear combination of traces. This is a generalization of \cite[Proposition~2.2]{BR} to the composite order case. Using this result, we give a counterexample to Borcherds's conjecture about vanishing of Tate cohomology. We propose a~weaker conjecture about vanishing of Tate cohomology.


% Original source lines 66--69
\section{Definitions and basic notions}\label{section2}
 The symbols $\mathbb{Z}$, $\mathbb{Q}$, $\mathbb{R}$, $\mathbb{C}$, $\mathbb{F}_p$ will be used to denote the integers, rational numbers, real numbers, complex numbers, the finite field of prime~$p$ order, respectively.

 In this section, we introduce some terminology in this paper and will give brief definitions.

\setcounter{Theorem}{4}
% Original source lines 137--179
\begin{Definition}
Let $G$ be a group and $N$ be an integer. An element $h$ of $G$ is $N$-regular if $(|h|,N)=1$.
\end{Definition}

\begin{Definition}
Let $R$ be a discrete valuation ring and $\mathfrak{m}_R$ the maximal ideal of $R$. Suppose that $k=R/\mathfrak{m}_R$ has prime characteristic $p$ and that the fraction field~$K$ of~$R$ has characteristic~$0$. The Brauer character of a $k[h]$-module~$M$ for a~$p$-regular element $h$ is defined by
\[
\widetilde{\operatorname{Tr}}(h|M)=\operatorname{Tr}\big(h|\widetilde{M}\otimes_RK\big),
\]
 where $\widetilde{M}$ is any $R$-free $R[h]$-module satisfying $M\cong\widetilde{M}\otimes_{R[h]}k[h]$.
\end{Definition}

\begin{Definition}[{\cite[Chapter IV]{CF}}]\label{Tate}
Suppose that $G$ is a group acting on an abelian group~$A$. Let ${\rm Nr}$ be the norm map $A\rightarrow A$; $a\mapsto\sum_{g\in G}{g(a)}$. Then, the $i$-th Tate cohomology ${\hat{H}}^i(G,A)$ is defined to be the following:
 \[
{\hat{H}}^i(G,A)= \begin{cases}
H_{-i-1}(G,A), & i\leq-2, \\
\operatorname{Ker}({\rm Nr})/\langle ga-a\,|\,a \in A, \, g\in G\rangle, & i=-1, \\
A^G/\operatorname{Im}({\rm Nr}), & i=0, \\
H^i(G,A), & i \geq1,
 \end{cases}
 \]
where $H_i(G,A)$ is the $i$-th homology, $H^i(G,A)$ is the $i$-th cohomology and $A^G$ is the largest submodule of $A$ on which $G$ acts trivially.
\end{Definition}

It has the following properties:
\begin{enumerate}\itemsep=0pt
\item[(1)] If $0\rightarrow A \rightarrow B \rightarrow C \rightarrow 0$ is a short exact sequence of $G$-modules, then we have a long exact sequence
\[
\cdots \rightarrow {\hat{H}}^{i-1}(G,C)\rightarrow {\hat{H}}^{i}(G,A)\rightarrow {\hat{H}}^{i}(G,B)\rightarrow {\hat{H}}^{i}(G,C)\rightarrow {\hat{H}}^{i+1}(G,A) \rightarrow \cdots.
\]
\item[(2)] $\bigoplus_{i\in \mathbb{Z}}{\hat{H}}^i(G,A)$ is $ |G |$-torsion module. In particular, if multiplication by $|G|$ is an isomorphism on $A$, then $\hat{H}^i(G,A)$ vanishes for all $i\in \mathbb{Z}$.
\end{enumerate}

Moreover, if $G$ is a cyclic group generated by $g$, then we have the following properties, where we denoted the Tate cohomology $\hat{H}^i(\langle g\rangle,A)$ for a cyclic group $\langle g\rangle$ by $\hat{H}^i(g,A)$.
\begin{enumerate}\itemsep=0pt
\item[(3)] ${\hat{H}}^i(g,A)\cong{\hat{H}}^{i+2}(g,A)$ for any $i\in\mathbb{Z}$.
\item[(4)]\label{Tsuper} ${\hat{H}}^\ast (g,A )=\hat{H}^0(g,A)\oplus\hat{H}^1(g,A)$ has more or less a super version of any $G$-invariant algebraic structure on~$A$. In particular, if~$V$ is a vertex algebra, then $\hat{H}^*(g,V)$ is a vertex superalgebra by the composite map
\[
\hat{H}^*(g,V)\otimes \hat{H}^*(g,V)\rightarrow \hat{H}^*(g,V\otimes V)\rightarrow \hat{H}^*(g,V((z)))\rightarrow \hat{H}^*(g,V)((z)).
\]
\item[(5)] If $A$ is a finite module, then $\big|{\hat{H}}^0(g,A)\big|=\big|{\hat{H}}^1(g,A)\big|$.
\end{enumerate}

\setcounter{Theorem}{9}
% Original source lines 193--210
\begin{Definition}
Let $R$ be a discrete valuation ring of mixed characteristic $(0,p)$, i.e., $R$ has chracteristic~0 and the residue field $R/\mathfrak{m}_R$ of $R$ has characteristic~$p$. Let~$S$ be the ring of integers of a~finite extension of the fraction field $\operatorname{Frac}(R)$ of $R$. Then $\mathfrak{m}^e_S=\mathfrak{m}_RS$ for some positive integer~$e$, called the ramification index. When $e=1$, $S$ is called an unramified extension of~$R$. On the other hand, when $e\neq 1$, $S$ is called a ramified extension of~$R$.
\end{Definition}

\begin{Definition}
Let $R$ be a commutative ring with prime characteristic~$p$. The Frobenius endomorphism $\operatorname{Frob}_p\colon R\rightarrow R$ is defined by $\operatorname{Frob}_p(r)=r^p$ for all $r\in R$.
\end{Definition}

If $R$ is a finite field, the Frobenius endomorphism is an automorphism.

\begin{Definition}
Let $A$ be a finite abelian group, $A^*$ be the homomorphisms $\operatorname{Hom}(A,\mathbb{C}^\times)$ from~$A$ to~$\mathbb{C}^\times$ and $f\colon A \rightarrow \mathbb{C}$ be a function. The discrete Fourier transform $\hat{f}\colon A^*\rightarrow \mathbb{C}$ is defined by $\hat{f}(\chi)=\sum_{a \in A}f(a)\chi (a)$ for any~$\chi \in A^*$.
\end{Definition}

There is an inverse discrete Fourier transform, it is described by the following
\[
f(a)=\frac{1}{|A|}\sum_{\chi \in A^*}\hat{f}(\chi)\overline{\chi}(a).
\]

\setcounter{section}{3}
% Original source lines 320--347
\section{Tate cohomology and ramified extensions}\label{section4}

In this section, we will prove properties of Tate cohomology related to ramified extensions and complete the proof of Theorem~4.7 of~\cite{BR} for an element of the monster of type~$15A $ or~$21A$.

Let $R$ be a discrete valuation ring. Let $g \in G$ be an element of order~$N$, $h \in G$ be an $N$-regular element and~$A$ be an $R[g]$-module.

\begin{Proposition}\label{tensor}
If a commutative ring homomorphism $R \rightarrow S$ is flat, then $\hat{H}^*(g,A\otimes_{R}S)\cong \hat{H}^*(g,A)\otimes_{R}S$.
\end{Proposition}

\begin{proof}Let ${\rm Nr}$ be a norm map as in Definition~\ref{Tate}. Then, $\hat{H}^0(g,A)\cong \operatorname{Ker}(g-1)/\operatorname{Im}({\rm Nr}).$ We have an exact sequence of $R$-modules $0\rightarrow \operatorname{Im}({\rm Nr})\rightarrow \operatorname{Ker}(g-1)\rightarrow \hat{H}^0(g,A)\rightarrow 0$. If $R \rightarrow S$ is flat, the sequence of $S$-modules $0\rightarrow \operatorname{Im}({\rm Nr})\otimes_R S\rightarrow \operatorname{Ker}(g-1)\otimes_R S\rightarrow \hat{H}^0(g,A)\otimes_R S\rightarrow 0$ is exact. We consider the following diagram:
\[ \xymatrix{
0\ar[r] & \operatorname{Im}({\rm Nr}\otimes_R {\rm id}_S)\ar[d] \ar[r] & \operatorname{Ker}((g-1)\otimes_R {\rm id}_S)\ar[d] \ar[r] & \hat{H}^0(g,A\otimes_R S)\ar[d]^{\psi_0} \ar[r] & 0 \\
 0\ar[r] & \operatorname{Im}({\rm Nr})\otimes_R S\ar[r] & \operatorname{Ker}(g-1)\otimes_R S\ar[r] & \hat{H}^0(g,A)\otimes_R S\ar[r] & 0.
 }
\]
Then, we have $\operatorname{Im}({\rm Nr}\otimes_R {\rm id}_S)\cong \operatorname{Im}({\rm Nr})\otimes_R S$ and $\operatorname{Ker}((g-1)\otimes_R {\rm id}_S)\cong \operatorname{Ker}(g-1)\otimes_R S$. Indeed, we can take the isomorphisms ${\rm Nr}\otimes_R {\rm id}_S(x\otimes s) \mapsto {\rm Nr}(x)\otimes s$ and $x\otimes s \mapsto x\otimes s$, respectively.
Also, the above diagram is commutative. Hence, by the Five lemma, $\psi_0$ is an isomorphism.

Like the degree 0 case, we have an exact sequence $0\rightarrow \operatorname{Im}(g-1)\rightarrow \operatorname{Ker}({\rm Nr})\rightarrow \hat{H}^1(g,A)\rightarrow 0$. As well above, we consider the following commutative diagram:
\[ \xymatrix{
0\ar[r] & \operatorname{Im}((g-1)\otimes_R {\rm id}_S)\ar[d] \ar[r] & \operatorname{Ker}({\rm Nr}\otimes_R {\rm id}_S)\ar[d] \ar[r] & \hat{H}^1(g,A\otimes_R S)\ar[d]^{\psi_1} \ar[r] & 0 \\
 0\ar[r] & \operatorname{Im}(g-1)\otimes_R S\ar[r] & \operatorname{Ker}({\rm Nr})\otimes_R S\ar[r] & \hat{H}^1(g,A)\otimes_R S\ar[r] & 0.
 }
\]
Then, we have $\operatorname{Im}((g-1)\otimes_R {\rm id}_S)\cong \operatorname{Im}(g-1)\otimes_R S$ and $\operatorname{Ker}({\rm Nr}\otimes_R {\rm id}_S)\cong \operatorname{Ker}({\rm Nr})\otimes_R S$. Indeed, we can take the isomorphisms $(g-1)\otimes_R {\rm id}_S(x\otimes s) \mapsto (gx-x)\otimes s$ and $x\otimes s \mapsto x\otimes s$, respectively.
By the Five lemma, $\psi_1$ is an isomorphism.
\end{proof}

\setcounter{Theorem}{4}
% Original source lines 391--420
Let $p$ be a prime factor of $N$ and $\zeta_N$ be a primitive $N$-th root of unity. Now, we consider a~ramified extension $R_p=\mathbb{Z}_p[\zeta_{N|h|}]$ of $\mathbb{Z}_p$.

We define rank 1 $R$-free $R_p[\langle g,h \rangle]$-modules $R_{n,m}$ with distinguished basis vectors $v_{n,m}\in R_{n,m}$, i.e.,
\[
R_{n,m}=\big\{cv_{n,m}\,|\,c\in R_p,\, gv_{n,m}=\zeta_N^nv_{n,m},\, hv_{n,m}=\zeta_{|h|}^mv_{n,m}\big\}.
\]
 Any $R_p[\langle g,h \rangle]$-module that is $R_p$-free of finite rank is an extension of rank~1 modules of the form~$R_{n,m}$ for $0\leq n\leq N-1$, $0\leq m\leq |h|-1$.
Hence, it suffices to consider the Tate cohomology of~$R_{n,m}$.

\begin{Lemma}\label{Tcoho}
For any $0\leq n \leq N-1$, $0\leq m \leq |h|-1$,
\begin{gather*}
\hat{H}^0(g,R_{n,m})= \begin{cases}
R_p/NR_p, & n=0, \\
0, & n\neq 0,
\end{cases}\\
\hat{H}^1(g,R_{n,m})= \begin{cases}
0, & n=0, \\
R_p/(1-\zeta_N^n)R_p, & n\neq 0,
\end{cases}
\end{gather*}
where $h$ acts as $\zeta_{|h|}^m$ on $R_p/NR_p$ and $R_p/\big(1-\zeta_N^n\big)R_p$.
\end{Lemma}

\begin{proof}
$R_{0,m}^g=R_{0,m}$, $R_{n,m}^g=0$, $n\neq 0$. Hence, $\hat{H}^0(g,R_{0,m})=R_p/NR_p$,
 $\hat{H}^0(g,R_{n,m})=0$, $n\neq 0$, because the norm map $R_{n,m} \rightarrow R_{n,m}$ is $v_{n,m} \mapsto \sum_{i=0}^{N-1}g^iv_{n,m}$.

By $\sum_{i=0}^{N-1}g^iv_{0,m}=Nv_{0,m}$, the kernel of the norm map on $R_{0,m}$ is~0. Hence, $\hat{H}^1(g,R_{0,m})=0$. By $\sum_{i=0}^{N-1}g^iv_{n,m}=\sum_{i=0}^{N-1}\zeta_N^{ni}v_{n,m}=0$, $n \neq 0$, the kernel of the norm map on~$R_{n,m}$ is~$R_{n,m}$. Therefore, by the definition of Tate cohomology, $\hat{H}^1(g,R_{n,m})=R_p/\big(1-\zeta_N^n\big)R_p$, $n\neq 0$.
\end{proof}


% Original source lines 422--597
\section[p-Brauer character]{$\boldsymbol{p}$-Brauer character}\label{section5}

In this section, we will define a new notion of generalized Brauer character and prove properties of it. We will give a counterexample to a conjecture of Borcherds about vanishing of Tate cohomology.

Let $p$ be a prime number. Let $R$ be a discrete valuation ring of mixed characteristic $(0,p)$, $v\colon R\rightarrow \mathbb{Z}_{\geq 0}\cup \{ \infty \}$ be a~surjective valuation on $R$ and $K=\operatorname{Frac}(R)$.

We can't use the Brauer character for $R_p/(1-\zeta_N^n)R_p$ because this is not a vector space over a field. To solve this, we will define a $p$-Brauer character.

\begin{Definition}\label{def5.1}
Let $G$ be a finite group and $M$ be a finite length $R[G]$-module. Let $h$ be a $p$-regular element of $G$. Suppose that $M$ has a composition series $0=M_0\subsetneq M_1\subsetneq \dots \subsetneq M_n=M$ as an $R[h]$-module.
 We define the $p$-Brauer character of $M$ for $h$ by
\[
\widetilde{\operatorname{Tr}}_p(h|M)=\cfrac{1}{v(p)} \sum_{i=1}^{n}\operatorname{Tr} \big(h|\widetilde{M_i}\otimes K\big),
\]
 where $\widetilde{M_i}$ is an $R$-free $R[h]$-module satisfying $M_i/M_{i-1} \cong\widetilde{{M}_i}\otimes_{R[h]}R/\mathfrak{m}_R[h]$.
\end{Definition}

\begin{Lemma}
This definition does not depend on our choices of $M_i$ and $\widetilde{M}_i$.
\end{Lemma}

\begin{proof}
For $M_i/M_{i-1} \cong\widetilde{{M}_i}\otimes_{R[h]}R/\mathfrak{m}_R[h]\cong\widetilde{{M}_i'}\otimes_{R[h]}R/\mathfrak{m}_R[h]$, we have $\operatorname{Tr}\big(h|\widetilde{M}_i\otimes K\big)=\widetilde{\operatorname{Tr}}(h|M_i/M_{i-1})=\operatorname{Tr}\big(h|\widetilde{M}_i'\otimes K\big)$. By the Jordan--H\"older theorem, the composition series of~$M$ is unique up to permutation and isomorphism. That is, for any two composition series $0=M_0\subsetneq M_1\subsetneq \dots \subsetneq M_n=M$, $0=M_0'\subsetneq M_1'\subsetneq \dots \subsetneq M_{n'}'=M$, we have $n=n'$ and some permutation $\sigma$ of $\{1,\dots,n\}$ s.t.\ $M_i/M_{i-1}\cong M_{\sigma(i)}'/M_{\sigma(i)-1}'$. Hence, $\sum_{i=1}^{n}\widetilde{\operatorname{Tr}}(h|M_i/M_{i-1})=\sum_{i=1}^{n}\widetilde{\operatorname{Tr}}(h|M_{\sigma(i)}'/M_{\sigma(i)-1}') =\sum_{i=1}^n\widetilde{\operatorname{Tr}}(h|M_i'/M_{i-1}')$. Therefore the $p$-Brauer character is well-defined.
\end{proof}

Note that it does depend on a choice of a prime number~$p$.

We can consider the notion of an irreducible $p$-Brauer character, so that $M$ is a simple $R[h]$-module that $M\otimes_{R[h]} R/\mathfrak{m}_R[h]$ is a simple $R/\mathfrak{m}_R[h]$-module. In particular, an irreducible $p$-Brauer is an irreducible Brauer character and a character table of $p$-Brauer characters is the same as of Brauer characters.

We will prove that the $p$-Brauer character is additive, where ``additive'' means that for any short exact sequence $0\rightarrow M\rightarrow M'\rightarrow M''\rightarrow 0$, $\widetilde{\operatorname{Tr}}_p$ satisfies $\widetilde{\operatorname{Tr}}_p(h|M')=\widetilde{\operatorname{Tr}}_p(h|M)+\widetilde{\operatorname{Tr}}_p(h|M'')$.

\begin{Lemma}\label{additive}
The $p$-Brauer character is additive on short exact sequences.
\end{Lemma}

\begin{proof}
Let $0\rightarrow M\xrightarrow{\varphi} M'\xrightarrow{\psi} M''\rightarrow 0$ be an exact sequence of finite length $R[h]$-modules. Suppose that $M$ has length $m$ and a composition series $0=M_0\subsetneq M_1\subsetneq \dots \subsetneq M_m=M$. We use analogous notation for~$M'$ and~$M''$. Then the length~$m'$ of~$M'$ is $m+m''$ and
\[
0=\varphi(M_0)\subsetneq \varphi(M_1)\subsetneq \dots \subsetneq \varphi(M_m)=\psi^{-1}(M''_0)\subsetneq \psi^{-1}(M''_1)\subsetneq \dots \subsetneq \psi^{-1}(M''_{m''})=M'
\]
 is a composition series of $M'$. By injectivity of $\varphi$, $M_i\cong \varphi(M_i)$ for any $i \in \{0, \dots,m\}$ and by surjectivity of $\psi$, $\psi^{-1}(M''_j)/\operatorname{Ker}(\psi)\cong M''_j$ for any $j \in \{0, \dots,m''\}$. Hence, we have $\varphi(M_i)/\varphi(M_{i-1})\cong M_i/M_{i-1}$ and
\[
\psi^{-1}(M''_j)/\psi^{-1}(M''_{j-1})\cong \big(\psi^{-1}(M''_j)/\operatorname{Ker}(\psi)\big)/\big(\psi^{-1}(M''_{j-1})/\operatorname{Ker}(\psi)\big) \cong M''_j/M''_{j-1}.
\]
 Therefore,
\begin{align*}
\widetilde{\operatorname{Tr}}_p(h|M')&=\frac{1}{v(p)}\sum_{i=1}^m\widetilde{\operatorname{Tr}}\big(h|\varphi(M_i)/\varphi(M_{i-1})\big) +\frac{1}{v(p)}\sum_{j=1}^{m''}\widetilde{\operatorname{Tr}}\big(h|\psi^{-1}(M''_j)/\psi^{-1}(M''_{j-1})\big) \\
 &=\frac{1}{v(p)}\sum_{i=1}^m\widetilde{\operatorname{Tr}}(h|M_i/M_{i-1})+\frac{1}{v(p)}\sum_{j=1}^{m''}\widetilde{\operatorname{Tr}}(h|M''_j/M''_{j-1}) \\
 &=\widetilde{\operatorname{Tr}}_p(h|M)+\widetilde{\operatorname{Tr}}_p(h|M'').\tag*{\qed}
\end{align*}\renewcommand{\qed}{}
\end{proof}

\begin{Lemma}\label{tensorBch}
Let $S$ be the ring of integers of a finite extension of $\operatorname{Frac}(R)$ and $A$ be a finite length $R[h]$-module. Then $\widetilde{\operatorname{Tr}}_p(h|A\otimes_{R}S)=\widetilde{\operatorname{Tr}}_p(h|A)$.
\end{Lemma}

\begin{proof}
We have that the $p$-Brauer character of $A$ is described the sum of $p$-Brauer characters of composition factors of $A$ by Definition~\ref{def5.1}. Hence, it suffices to consider the case $A$ is a simple $R[h]$-module on which $h$ acts as $\zeta_{|h|}^m$. Now,
\[
\widetilde{\operatorname{Tr}}_p(h|A)=\frac{1}{v_R(p)}\operatorname{Tr} \big(h|\widetilde{A}\otimes \operatorname{Frac}(R)\big),
\]
where $v_R$ is a discrete valuation on~$R$, $\widetilde{A}$ is an $R$-free $R[h]$-module satisfying $A\cong \widetilde{A}\otimes_{R[h]} R/\mathfrak{m}_R[h]$ and $\mathfrak{m}_R$ is a maximal ideal of~$R$. We have
\[
A\otimes_R S\cong \widetilde{A}\otimes_{R[h]} R/\mathfrak{m}_R[h]\otimes_R S\cong \widetilde{A}\otimes_{R[h]} S/\mathfrak{m}_RS[h].
\]

 Let $v_S$ be a discrete valuation on $S$ and $\mathfrak{m}_S$ be a maximal ideal of~$S$. If $S$ is an unramified extension, then $\mathfrak{m}_S=\mathfrak{m}_RS$ and $v_S(p)=v_R(p)$. Hence,
 \begin{align*}
\widetilde{\operatorname{Tr}}_p(h|A\otimes_R S) &=\frac{1}{v_S(p)}\operatorname{Tr}\big(h|\widetilde{A}\otimes_R S\otimes \operatorname{Frac}(S)\big)\\
& =\frac{1}{v_R(p)}\operatorname{Tr}\big(h|\widetilde{A}\otimes \operatorname{Frac}(R)\big)=\widetilde{\operatorname{Tr}}_p(h|A).
 \end{align*}

 If $S$ is a ramified extension, then there is a positive integer $k>1$ s.t.\ $\mathfrak{m}_S^k=\mathfrak{m}_RS$ and $v_S(p)=kv_R(p)$. Hence, we have $A\otimes_R S=\widetilde{A}\otimes_{R[h]}S/\mathfrak{m}_S^k[h]$ and a composition series $0\subsetneq \widetilde{A}\otimes S/\mathfrak{m}_S\subsetneq \dots \subsetneq \widetilde{A}\otimes S/\mathfrak{m}_S^k$. Its composition factors are all isomorphic to $\widetilde{A}\otimes_{R[h]} S/\mathfrak{m}_S[h]$. Therefore,
\begin{align*}
\widetilde{\operatorname{Tr}}_p(h|A\otimes_R S) & =\frac{1}{v_S(p)}\sum_{i=1}^k\operatorname{Tr} \big(h|\widetilde{A}\otimes_R S\otimes \operatorname{Frac}(S)\big)\\
& =\frac{1}{kv_R(p)}k\operatorname{Tr}\big(h|\widetilde{A}\otimes \operatorname{Frac}(R)\big) =\widetilde{\operatorname{Tr}}_p(h|A).\tag*{\qed}
\end{align*}\renewcommand{\qed}{}
\end{proof}

By Proposition \ref{tensor} and Lemma~\ref{tensorBch}, any finitely generated $\mathbb{Z}$-free $\mathbb{Z}[\langle g,h \rangle]$-module base changes to an $R_p[\langle g,h \rangle]$-module that is $R_p$-free of finite rank and the $p$-Brauer character of the Tate cohomology is unchanged.
By Lemma \ref{additive}, it suffices to consider the $p$-Brauer character for the Tate cohomology of~$R_{n,m}$.

We will calculate the Tate cohomology of $R_{n,m}$ by using the $p$-Brauer character. To do this, we use the following lemma:

\begin{Lemma}[{\cite[Chapter III, Lemma 3]{CF}}] \label{pR}
 Let $q=p^k$ for a positive integer~$k$. Then $pR=(1-\zeta_q)^{\varphi(q)}R$, where $\varphi$ is Euler's totient function.
\end{Lemma}

For a positive integer $s$ and $R=\mathbb{Z}_{p}[\zeta_{p^s}]$, a discrete valuation $v(p)$ is given by $\varphi(p^s)$ and maximal ideal of~$R$ is $(1-\zeta_{p^s})R$.

\begin{Lemma}\label{Bch}
Suppose that the prime factorization of $N$ is $\prod p_i^{n_i}$.
Then,
\[
\widetilde{\operatorname{Tr}}_{p_i}(h|\hat{H}^*(g,R_{n,m}))= \begin{cases}
n_i\zeta_{|h|}^m, & n=0, \vspace{1mm}\\
\dfrac{-\zeta_{|h|}^m}{\varphi \big(p_i^l\big)}, & |g^n|=p_i^l, \ 1\leq l\leq n_i, \vspace{1mm}\\
0, & otherwise.
\end{cases}.
\]
\end{Lemma}

\begin{proof}Now, $h$ acts as $\zeta_{|h|}^m$ on $\hat{H}^*(g,R_{n,m})$ and $v(p)=\varphi\big(p_i^{n_i}\big)$. By \[
\widetilde{\operatorname{Tr}}_{p_i}\big(h|\hat{H}^*(g,R_{n,m})\big)=\widetilde{\operatorname{Tr}}_{p_i}\big(h|\hat{H}^0(g,R_{n,m})\big) -\widetilde{\operatorname{Tr}}_{p_i}\big(h|\hat{H}^1(g,R_{n,m})\big)
\]
 and Lemma~\ref{Tcoho}, it suffices to calculate the $p_i$-Brauer characters of $R_{p_i}/NR_{p_i}$ and $R_{p_i}/(1-\zeta_N^n)R_{p_i}$.

 When $n=0$, by Lemmas~\ref{Tcoho} and~\ref{pR},
\[
 R_{p_i}/NR_{p_i}\cong R_{p_i}/p_i^{n_i}R_{p_i}\cong R_{p_i}/\big(1-\zeta_{p_i^{n_i}}\big)^{n_i\varphi(p_i^{n_i})}R_{p_i}.
\]
This module has a composition series
\[
0\subsetneq R_{p_i}/(1-\zeta_{{p_i}^{n_i}})R_{p_i}\subsetneq \dots \subsetneq R_{p_i}/\big(1-\zeta_{{p_i}^{n_i}}\big)^{n_i\varphi(p_i^{n_i})}R_{p_i}
\]
 and its composition factors are all isomorphic to $R_{p_i}/\big(1-\zeta_{{p_i}^{n_i}}\big)R_{p_i}$. Hence,
\[
 \widetilde{\operatorname{Tr}}_{p_i}(h|R_{p_i}/NR_{p_i})=n_i\zeta_{|h|}^m.
\]

 When $|g^n|=p_i^l$, $\zeta_N^n$ is a $p_i^l$-th primitive root of unity. By Lemma~\ref{pR}, we have
\[
 \big(1-\zeta_{p_i^l}\big)^{\varphi(p_i^l)}=\big(1-\zeta_{p_i^{n_i}}\big)^{\varphi(p_i^{n_i})}.
\]
 Hence, by Lemma~\ref{Tcoho}, we have
\[
R_{p_i}/\big(1-\zeta_{{p_i}^l}\big)R_{p_i}\cong R_{p_i}/\big(1-\zeta_{{p_i}^{n_i}}\big)^{\varphi(p_i^{n_i})/\varphi(p_i^l)}R_{p_i}.
\]
 This module has the length $\varphi\big(p_i^{n_i}\big)/\varphi\big(p_i^l\big)$ and its composition factors are all isomorphic to $R_{p_i}/\big(1-\zeta_{{p_i}^{n_i}}\big)R_{p_i}$. Calculating the $p_i$-Brauer character, we have
\[
 \widetilde{\operatorname{Tr}}_{p_i}\big(h|R_{p_i}/\big(1-\zeta_{p_i^l}\big)R_{p_i}\big)=1/\varphi\big(p_i^l\big).
\]

 When $|g^n|=p_i^kq$, where $q\neq 1$, $(q,p_i)=1$ and $0\leq k\leq n_i$, $\zeta_N^n$ is a $p_i^kq$-th primitive root of unity. Let $\pi$ be a surjective homomorphism $R_{p_i}\rightarrow \mathbb{F}_{p_i}\big[\zeta_{N|h|/p_i^{n_i}}\big]$; $\sum_{j=0}^\infty c_jp_i^j \mapsto c_0$, $\zeta_{N|h|} \mapsto \zeta_{N|h|/p_i^{n_i}}$. Then $\operatorname{Ker}(\pi)$ is the maximal ideal of $R_{p_i}$ and $\operatorname{Frob}_{p_i}$ on $\mathbb{F}_{p_i}[\zeta_{N|h|/p_i^{n_i}}]$ is an automorphism. We have $\pi\big(\zeta_{p_i^kq}\big)=\operatorname{Frob}_{p_i}^{-k}\big(\zeta_{p_i^kq}^{p_i^k}\big)\neq 1$.
 Hence, $1-\zeta_N^n \notin \operatorname{Ker}(\pi)$. $1-\zeta_N^n$ is a unit of $R_{p_i}$ because every element of a local ring not included in a maximal ideal is a unit. Therefore, $R_p/\big(1-\zeta_N^n\big)R_p=0$.
\end{proof}

We generalize Proposition~\ref{prop2.2} \cite[Proposition~2.2]{BR} using the $p$-Brauer character.

\begin{Theorem} \label{2.2kai}
Let $G$ be a finite group. Suppose that $g \in \operatorname{Cent}(G)$ has order $N$ and that $h \in G$ is an $N$-regular element. Let~$A$ be a finitely generated $R$-free $R[\langle g,h \rangle]$-module. Then $\hat{H}^*(g,A)=\hat{H}^0(g,A)-\hat{H}^1(g,A)$ is a virtual representation of $\langle h \rangle$, and
\begin{gather*}
\widetilde{\operatorname{Tr}}_{p_i}(h|\hat{H}^*(g,A))=\sum_{k=1}^{N-1}a_{k,p_i}\operatorname{Tr}\big(g^kh|A\big) ,
\\
a_{k,p_i}=\begin{cases}
\dfrac{1}{\sum_{(p_i,d)=1,d|N}\varphi (N/d)}\left(n_i-l-\dfrac{n_i-l-1}{p_i}\right), & \big(k,p_i^{n_i}\big)=p_i^l, \ 0\leq l \leq n_i-1, \\
0, & \big(k,p_i^{n_i}\big)=p_i^{n_i},
\end{cases}
\end{gather*}
where the prime factorization of $N$ is $\prod p_i^{n_i}$.
\end{Theorem}

\begin{proof}
Suppose
\[
\widetilde{\operatorname{Tr}}_{p_i}\big(h|\hat{H}^*(g,A)\big)=\sum_{k=1}^{N-1}a_{k,p_i}\operatorname{Tr}\big(g^kh|A\big). \]
It suffices to take $A=R_{n,m}$.
By an inverse discrete Fourier transform and Lemma~\ref{Bch}, we have
\begin{align*}
a_{k,p_i}&=\frac{1}{N}\sum_{b=0}^{N-1}\widetilde{\operatorname{Tr}}_{p_i}\big(h|\hat{H}^*(g,R_{b,m})\big)\overline{\operatorname{Tr}\big(g^kh|R_{b,m}\big)} \\
 &=\frac{1}{N}\Bigg(n_i+\sum_{|g^b|=p_i} \frac{-1}{\varphi (p_i)}\zeta_N^{-kb}+\sum_{ |g^b|=p_i^2} \frac{-1}{\varphi \big(p_i^2\big)}\zeta_N^{-kb}+\dots+\sum_{|g^b|=p_i^{n_i}} \frac{-1}{\varphi \big(p_i^{n_i}\big)}\zeta_N^{-kb}\Bigg).
\end{align*}

Let $\big(k,p_i^{n_i}\big)=p_i^l$. If $s>l$, then $\sum_{|g^b|=p_i^s}\zeta_N^{-kb}$ is a sum of primitive $p_i^{s-l}$-th roots of unity and equals an integer multiple of the M\"obius function $\mu\big(p_i^{s-l}\big)$. If $s\leq l$, then we have
\[
\sum_{|g^b|=p_i^s}\zeta_N^{-kb}=\varphi(p_i^s).
\]
Hence, when $\big(k,p_i^{n_i}\big)=p_i^{n_i}$, we have 
\[ a_{k,p_i}=\frac{1}{N}(n_i-1-\dots-1)=0.\]
When $\big(k,p_i^{n_i}\big)=p_i^{l}$, because $\big(N/p_i^{n_i}\big)=\sum_{d|(N/p_i^{n_i})}\varphi\big(N/p_i^{n_i}d\big)$, we have
\begin{align*}
a_{k,p_i}&=\frac{1}{N}\Bigg(n_i-1-\dots-1+\frac{p_i^l}{\varphi \big(p_i^{l+1}\big)}\Bigg)\\
 &=\frac{(n_i-l)\varphi \big(p_i^{l+1}\big)+p_i^l}{N\varphi \big(p_i^{l+1}\big)} =\frac{p_i^{l+1}(n_i-l)-p_i^l(n_i-l-1)}{p_i^{l+1}\varphi\big(p_i^{n_i}\big)\big(N/p_i^{n_i}\big)}\\
 &=\frac{1}{\sum_{(p_i,d)=1,d|N}\varphi (N/d)}\bigg(n_i-l-\frac{n_i-l-1}{p_i}\bigg).\tag*{\qed}
\end{align*}\renewcommand{\qed}{}
\end{proof}
\begin{thebibliography}{99}
\bibitem[3]{BR}
Borcherds R.E., Ryba A.J.E., Modular {M}oonshine.~{II}, \href{https://doi.org/10.1215/S0012-7094-96-08315-5}{\textit{Duke Math.~J.}}
 \textbf{83} (1996), 435--459.

\bibitem[5]{CF}
Cassels J.W.S., Fr\"ohlich A. (Editors), Algebraic number theory, Academic
 Press, London, Thompson Book Co., Inc., Washington, D.C., 1967.

\end{thebibliography}
\end{document}
